%!TEX root = conference_101719.tex


%Scaling model capacity has vastly improved instruction following~\cite{zhang2023instruction}, tool use~\cite{qin2023toolllm}, and reasoning abilities~\cite{xu2025largereasoningmodelssurvey} in many 

The staggering energy costs of both inference and training with large language models has made energy a growing concern in NLP~\cite{strubell19,fernandez-etal-2025-energy,huang2026eop}, Machine Learning~\cite{Jesse-dodge-icml,ICLR2025_8e016430}, and systems research~\cite{wilkins2024offlineenergyoptimalllmserving}. 
%Large Language Models~\footnote{In this work, we consider LLMs to be models with billions of parameters.} (LLMs) are incurring staggering energy costs for both inference and training. 
Recent estimates show that energy consumption for LLM inference alone accounts for 80-90\% of total energy consumption in data centers~\cite{Wenjen_2024}. Model developers and service providers today have a design space that spans models of various sizes, optimizations, and parallelization options with a wide-ranging energy costs. A key deployment question is how to make energy-efficient choices from this vast design space.


%%, and the inference energy is projected to grow to 1,050 terawatt-hours by 2026~\cite{kakolyris2024sloawaregpufrequencyscaling}. 

%Large Language Models~\footnote{In this work, we consider LLMs to be models with billions of parameters.} (LLMs) 
%including Llama, Olmo, GPT4, and Claude have vastly improved a wide variety of tasks including instruction following~\cite{zhang2023instruction}, tool use~\cite{qin2023toolllm}, and reasoning abilities~\cite{xu2025largereasoningmodelssurvey}. However, these  
%\emph{increase in compute and energy costs}.
%Coupled with the widespread adoption and deployments for web scale queries\footnote{OpenAI services have hit a billion queries per day as of April 2025~\cite{chatgpt-stats}.}, these advances also raise \be{LLM inference energy} to staggering levels. 
%Recent estimates show that energy consumption for LLM inference can account for 80--90\% of total energy consumption in certain data centers~\cite{Wenjen_2024}, and the inference energy alone is projected to grow to 1,050 terawatt-hours by 2026~\cite{kakolyris2024sloawaregpufrequencyscaling}.


%Consequently, model developers and service providers increasingly need to make energy-aware choices, since even a small change in design choice can result is substantial increase in energy cost. 

The main bottleneck is that service providers do not have the tools to easily characterize the LLM energy consumption of their LLM serving configuration. A natural approach is to measure inference energy directly using full-system hardware energy monitors~\cite{anthony2020carbontrackertrackingpredictingcarbon, energyVis, mauricio, benoit_courty_2024_11171501}. However, in multi-tenant data centers, hardware power monitors are either not available %since they require inline meters and administrative privileges 
or even if available, require exclusive machine access. 

Energy prediction offers a promising alternative to direct measurement and prior work has made progress in predicting the energy consumption of LLM inference~\cite{strubell19,Nguyen-hotcarbon,wilkins2024offlineenergyoptimalllmserving,irene,Jesse-dodge-icml,codecarbon}. However, all of these approaches assume that LLMs are served on a \emph{single GPU}. In practice, modern LLMs routinely employ parallel inference over multiple GPUs~\cite{sangjin23_envpipe,megatron}.

The \be{communication overheads introduced by multi-GPU execution fundamentally alter the energy profile}, making existing prediction models inaccurate (see \S\ref{s:challenges}). As a consequence, developers lack practical tools for answering crucial deployment questions, such as: should the model be parallelized across 2- or 4- GPUs for energy efficiency? Would pipeline or tensor parallelism provide better energy savings for a given model?

\noindent In this work, we present \textbf{Parallelized Inference Energy Predictor (\sys{})}, an accurate energy prediction framework for multi-GPU LLM inference that generalizes across model architectures. 
%\anshul{clarify whether PIE-P generalizes across model architectures, h/w, or both.}
\sys{} is designed to support energy-efficient parallelization decisions across data, pipeline, and tensor parallelism.
Beyond predicting the energy consumption of the entire model, \sys{} also estimates the energy of individual modules, enabling developers to identify energy bottlenecks during inference. 

To obtain accurate energy predictions in parallel inference deployments, 
\sys{} needs to keep track both parallel communication and computation which are both independently challenging. Parallel inference incurs substantial energy from inter-GPU communication and transfers (see \S\ref{s:challenges}),
which are absent in single-GPU settings and can account for a large fraction of end-to-end energy (see Figure~\ref{fig:all_reduce_energy}). 
%%In addition, multi-GPU execution introduces run-to-run variability from stragglers and synchronization (see \S\ref{}), creating non-deterministic energy consumption.
In terms of computation, energy depends how each module in the model is executed on the GPU. % which in-turn depends on the tensor shape and hardware/resource characteristics
In a multi-GPU setting, one has to account for how the modules are distributed across GPUs based on the parallelization strategy.




\begin{comment}
\begin{itemize}[leftmargin=*,itemsep=0pt]
\item \textbf{Communication energy.}
Parallel inference incurs substantial energy from inter-GPU collectives and transfers (see \S\ref{}),
which are absent in single-GPU settings and can account for a large fraction of end-to-end energy (see Figure~\ref{fig:all_reduce_energy}).
Ignoring this component leads to systematic underestimation of energy.

\item \textbf{Non-deterministic synchronization.}
Multi-GPU execution introduces run-to-run variability from stragglers and synchronization (see \S\ref{}), creating non-deterministic energy consumption. This
variability makes both measurement and prediction more challenging.
%variability makes both measurement (separating transfer vs.\ waiting) and prediction (capturing the distribution of synchronization energy rather than a single trace).}
\anshul{do we address this?}
\end{itemize}
\end{comment}

To address these challenges,
\sys{} directly captures the fine-grained energy use of compute and communication. % and accounts for non-determinism.
Specifically, we adopt the \emph{model-tree abstraction} (similar to existing approaches such as IrEne~\cite{irene, neuralpower}) to decompose the model energy into individual constituents. But different from existing decomposition techniques, \sys{} uses a \be{parallelism-aware model decomposition} that takes into account both (i) how the modules are distributed across the GPUs and (ii) adds inter-GPU communication events as individual modules in the model-tree abstraction.
%

To learn the energy use of these communication modules across different parallelisms, \sys{} identifies \emph{structural model features}, such as number of attention heads, to capture the relationship between model architecture and communication patterns.
% To account for parallelization-induced synchronization overheads, \sys{} \be{models inter-GPU collective latency} as a function of fixed communication latency, effective bandwidth, and the bytes sent/received.
% At runtime, \sys{} then employs this model to predict communication time and estimate the corresponding energy contribution for communication.
%
\sys{} \be{models communication energy} as calibrated communication power times predicted latency, which is parameterized by effective bandwidth and per-GPU communication volume.
The resulting communication-module estimates are then composed with compute-module estimates through a multi-level regressor to produce the final end-to-end energy prediction.
% \sys{} \be{models inter-GPU communication latency} as a function of the communication latency, effective bandwidth, and the bytes sent/received.
%The predicted communication latency is then used as part of the communication-module feature set, and
% \sys{} then composes the resulting communication-module estimates with compute-module estimates through a multi-level regressor to produce the final end-to-end energy prediction.

%To enable scalable energy predictions, \sys{} employs \emph{aggregate feature representation} to concisely track the features of multiple GPUs
%based on aggregates of their runtime features. 

\begin{comment}
To address these challenges,
\sys{} (i) introduces a \be{parallelism-aware model decomposition} that directly captures the fine-grained energy use of compute and communication (including collectives such as AllReduce), including accounting for non-determinism; and (ii) combines these components by building a \be{hierarchical, module-level regression model} to produce accurate end-to-end energy predictions for any parallelism configuration. 

As a first step, \sys{} adopts a \emph{model-tree} abstraction for energy prediction (similar to existing approaches such as IrEne~\cite{irene, neuralpower}). 
However, \sys{} goes beyond the single-GPU setting considered in these prior works. % and specifically overcomes the challenges in predicting energy consumption for parallel inference. 
To this end, \sys{} makes the following contributions: 
\begin{enumerate}
    \item 
\emph{Communication Energy Profiling} by treating inter-GPU communication events as individual modules in the model-tree abstraction and learning their energy use via benchmarking.
    \item \emph{Use of Structural Model Features}, such as number of attention heads, to capture the relationship between model architecture and communication patterns when using different parallelisms.
    \item {\emph{Communication-Aware Prediction} to account for parallelization-induced communication overheads by modeling inter-GPU collective latency as a function of fixed communication latency, effective bandwidth, and the bytes sent/received, and using the predicted communication time to estimate the corresponding AllReduce energy contribution.}
    \item \emph{Aggregate Runtime Feature Representation} to scalably and concisely represent the features of multiple GPUs
    based on aggregates of their runtime features. 
\end{enumerate}
\aruna{(For Anshul) There is some repetition between this set of contributions and the previous paragraph. I wonder if we can combine the two in any way.}
\end{comment}

% We implement \sys{} to predict the energy consumption on a variety of open LLM  families (Vicuna, Mistral, Llama, and Qwen) of varying sizes (from 1B to 70B \aruna{check}), on three different hardware (\aruna{name the hardware}) and three different parallelization strategies.   

We implement \sys{} to predict energy consumption across open LLM families
(Vicuna, Mistral, Llama, and Qwen), 
%with model sizes ranging from 7B to 70B parameters, 
on three GPU platforms (A6000, B40, H200), and across
%\youngwon{we mention which hardware and which models in Section 5 as well} We evaluate \sys{} across 
tensor, pipeline, and data
parallelism. % with tensor parallelism serving as the primary evaluation setting
%because it introduces the most challenging communication and synchronization overheads.
nder tensor-parallel evaluations, \sys{ achieves 15.5-17.6\% average model-level MAPE across A6000, B40, and H200. Average MAPE under pipeline and data parallelism is 12.2\% and 10.8\%, respectively.}

e compare \sys{ with CodeCarbon~\cite{codecarbon}, Wilkins et al.~\cite{wilkins2024offlineenergyoptimalllmserving}, IrEne~\cite{irene}, NVML-based regression~\cite{nvml}, NeuSight adapted to predict energy~\cite{neusight}, and IrEne augmented with historical communication-energy measurements.} 
%\aruna{We also compare \sys{} against Neusight~\cite{neusight}, which we modify to predict energy rather than its original latency objective, as well as an extended version of IrEne that accounts for network energy based on historical measurement data.} 
Compared to the next-best performing baseline, \sys{} typically reduces energy prediction error by 1.5-3$\times$.
We also show that \sys{} generalizes to unseen models and accurately predicts module-level energy consumption, and we conduct ablation studies to analyze its behavior.
We end with a use case that shows how \sys{} can help navigate the trade-off between inference time and energy consumption across different model sizes and GPU configurations. We commit to releasing all code and data upon publication.